% ============================================================
% Lecture-specific additions on top of the unchanged template
% (template-preamble.tex = chapter_00.tex lines 1-87, verbatim).
% Everything visual comes from the template: Madrid theme, accent
% colour, headline navigation, frame-number footer, tcolorbox styles.
% ============================================================

% --- Section outline frames -------------------------------------------------
% The template opens every section with an automatic "Outline" frame. A
% 60-minute lecture with a live build cannot spend six frames on outlines, so
% they are switched off here; the overview slide and the template's headline
% navigation take over that role.
\AtBeginSection[]{}

% --- Speaker notes ----------------------------------------------------------
% main.tex hides notes; main-notes.tex defines \shownotes for a two-screen
% edition (slide left, notes right), as in the source deck.
\ifdefined\shownotes
  \usepackage{pgfpages}
  \setbeameroption{show notes on second screen=right}
  \setbeamertemplate{note page}{%
    \vspace*{2mm}\hspace*{4mm}%
    \begin{minipage}{0.88\paperwidth}\footnotesize\raggedright\insertnote\end{minipage}}
\else
  \setbeameroption{hide notes}
\fi

% --- Bibliography --------------------------------------------------------------
% The source deck used biber. On the preparation machine biber fails to start
% ("extracting arm64 binary with lipo failed"), so this deck uses biblatex with
% the bibtex backend, which needs nothing beyond a standard TeX Live.
\usepackage{csquotes}
\usepackage[style=authoryear, backend=bibtex, maxbibnames=4, giveninits=true,
            uniquename=false, uniquelist=false, url=true, doi=true, isbn=false]{biblatex}
\addbibresource{references.bib}
\renewcommand*{\bibfont}{\fontsize{6.8}{8.1}\selectfont}
\setlength{\bibitemsep}{0.15ex}
\setbeamertemplate{bibliography item}{}
\AtEveryBibitem{\clearfield{urlyear}\clearfield{urlmonth}\clearfield{urlday}}

% --- Title-slide content (from the source deck's title slide) ---------------
\title{AI or Reality? Can We Tell the Difference?}
\subtitle{Fluent text, false claims and a live experiment with language models}
\author{Prof.\ Dr.\ Christoph Weisser}
\institute{Hochschule Bielefeld -- University of Applied Sciences and Arts (HSBI)}
\date{2026}
\graphicspath{{images/}}
\newcommand{\mailme}{\href{mailto:christoph.weisser@hsbi.de}{\color{accent}\texttt{christoph.weisser@hsbi.de}}}
\titlegraphic{\small\faEnvelope[regular]\ \mailme\\[1mm]
  {\scriptsize\color{black!55}All images are credited in the appendix.}}
\hypersetup{pdftitle={AI or Reality? Can We Tell the Difference?},
            pdfauthor={Prof. Dr. Christoph Weisser}, colorlinks=true,
            urlcolor=accent, linkcolor=accent, citecolor=accent}

% --- Course repositories ------------------------------------------------------
\newcommand{\repoPython}{\href{https://github.com/ChrisW09/Python-for-AI-Driven-Automation}{\color{accent}\faGithub\ \texttt{github.com/ChrisW09/Python-for-AI-Driven-Automation}}}
\newcommand{\repoQRM}{\href{https://github.com/ChrisW09/Quantitative-Research-Methods}{\color{accent}\faGithub\ \texttt{github.com/ChrisW09/Quantitative-Research-Methods}}}

% --- Small helpers in the template's style ------------------------------------
% Compact visible source line, same look as the template's \islpsource.
\newcommand{\slidesource}[1]{\vfill\hfill{\tiny\itshape\color{gray}#1}\par}
% A/B badge for the quiz rounds (from the source deck, recoloured to the accent).
\newcommand{\quizbadge}[2][accent]{\tikz[baseline=(b.base)]{\node[circle, fill=#1, text=white,
  font=\bfseries\large, minimum size=7mm, inner sep=0pt] (b) {#2};}}
% Token chips for the tokenisation slide (from the source deck, recoloured).
\newcommand{\tok}[1]{\tikz[baseline=(tk.base)]{\node[inner xsep=1.05mm, inner ysep=0.6mm,
  rounded corners=0.8mm, fill=accent!14, text=black] (tk) {\ttfamily\footnotesize #1};}\hspace{0.6mm}}
\newcommand{\tokcount}[2]{\begin{minipage}[c]{19mm}\raggedleft
  {\footnotesize\bfseries\color{accent}#1}\\[0.2mm]{\tiny\color{black!55}#2}\end{minipage}\hspace{2.2mm}}
% Inline step number for walkthroughs.
\newcommand{\step}[1]{\tikz[baseline=(s.base)]{\node[circle, fill=accent, text=white,
  font=\bfseries\scriptsize, minimum size=4.2mm, inner sep=0pt] (s) {#1};}\,}
% A keyboard-style label for exact UI names.
\newcommand{\ui}[1]{\tikz[baseline=(u.base)]{\node[draw=black!35, fill=black!4, rounded corners=0.7mm,
  inner xsep=1mm, inner ysep=0.5mm, font=\scriptsize\sffamily] (u) {#1};}}
% Prompt listings: plain text, no syntax highlighting, no shell escape.
\lstdefinestyle{prompt}{language={}, basicstyle=\ttfamily\scriptsize, columns=fullflexible,
  keepspaces=true, breaklines=false, showstringspaces=false, upquote=true,
  backgroundcolor=\color{black!3}, frame=leftline, rulecolor=\color{accent},
  framerule=1.2pt, xleftmargin=2mm, framexleftmargin=1mm, aboveskip=1mm, belowskip=0mm,
  keywordstyle={}, commentstyle={}, stringstyle={}, literate={}}
\lstdefinestyle{promptbig}{style=prompt, basicstyle=\ttfamily\footnotesize}
